\section{Ameaças à Validade}
\label{sec:threats}

Esta seção explicita os principais fatores que limitam o alcance das conclusões e define como os resultados devem ser interpretados.
O objetivo não é enfraquecer o estudo, mas delimitar com precisão o que os dados sustentam, em que condições a comparação é válida e quais extensões ainda são necessárias.

\subsection{Validade Interna}

A primeira ameaça à validade interna decorre da heterogeneidade de regime entre os dois blocos analíticos do artigo.
Os métodos clássicos foram avaliados em cenário multissensor, com forte preocupação com viabilidade computacional e equilíbrio entre sensores e níveis de ruído.
Os modelos de aprendizado profundo, por sua vez, foram avaliados no domínio Sentinel-2, coerente com seu regime de treinamento.
Essa diferença não inviabiliza o estudo, mas impede interpretar a presença simultânea dos dois blocos como evidência de um ranking universal entre todas as técnicas.
Por isso, as conclusões foram formuladas como comparações internas a cada família metodológica.

Uma segunda ameaça interna está associada ao custo computacional muito desigual entre os métodos clássicos.
Para preservar comparabilidade prática, a avaliação consolidada desse bloco foi realizada sobre um subconjunto balanceado de recortes.
Essa decisão é metodologicamente defensável, porque todos os métodos clássicos foram submetidos exatamente ao mesmo protocolo, mas reduz a cobertura empírica em relação ao universo completo de amostras disponíveis no repositório.
Assim, a inferência principal deve ser lida como comparação relativa sob desenho controlado, e não como medição exaustiva de desempenho em toda a base potencial.

A terceira ameaça interna diz respeito à variabilidade experimental efetivamente observada nos resultados agregados.
Embora o pipeline contemple controle de sementes e reamostragem, parte importante das tabelas finais reflete variação entre recortes, sensores e níveis de ruído mais do que replicações independentes de treinamento ou execução para todas as combinações possíveis.
Isso significa que os intervalos de confiança resumem incerteza empírica relevante, mas não eliminam por completo o risco de sensibilidade a escolhas específicas de amostragem, particionamento e configuração.

\subsection{Validade de Construção}

Uma ameaça de construção vem da própria escolha das métricas.
PSNR, SSIM, RMSE, SAM e ERGAS capturam aspectos centrais da reconstrução, mas não são equivalentes a utilidade em tarefa final.
Um método pode apresentar ótimo desempenho radiométrico e espectral sem necessariamente produzir a melhor entrada para classificação, detecção de mudança ou estimativa biofísica.
Portanto, este artigo avalia a qualidade da reconstrução em si, e não o impacto downstream em toda cadeia de sensoriamento remoto.

Outra limitação de construção envolve a interpretação da entropia local.
Ela foi utilizada como variável organizadora da análise por representar complexidade espacial da cena, mas não esgota todos os fatores que afetam a dificuldade do preenchimento de lacunas.
Geometria de bordas, textura direcional, autocorrelação espacial, presença de objetos raros e contraste espectral entre bandas também podem influenciar fortemente o comportamento dos métodos.
Logo, entropia não deve ser lida como sinônimo completo de dificuldade, e sim como um moderador útil, porém parcial, do regime de erro.

Há ainda uma limitação associada à discretização por tercis e ao uso de janelas fixas de entropia.
Essa escolha facilita comparação entre tabelas e figuras, mas comprime uma estrutura contínua em faixas agregadas.
Com isso, transições locais mais finas podem ser atenuadas, especialmente em cenas cuja complexidade espacial muda de modo gradual no entorno da lacuna.

\subsection{Validade Estatística}

Do ponto de vista estatístico, a principal ameaça está na dependência entre variáveis explicativas derivadas em múltiplas escalas.
As entropias calculadas em janelas 7x7, 15x15 e 31x31 são naturalmente correlacionadas entre si.
Por essa razão, modelos multivariados com todas as escalas simultaneamente são úteis para detectar associação global entre estrutura da cena e desempenho, mas não devem ser interpretados de maneira causal ou isolada coeficiente a coeficiente.

Também é necessário cautela ao interpretar significância estatística em um conjunto amplo de comparações.
O uso de correção por \gls{FDR}, testes não paramétricos e bootstrap em cluster reduz o risco de inferências espúrias, mas não substitui análise substantiva do tamanho de efeito.
Diferenças pequenas podem atingir significância em cenários com muitas observações, enquanto diferenças operacionalmente importantes podem aparecer de forma mais localizada e depender do regime de ruído ou da faixa de entropia.

Além disso, a reamostragem por \texttt{patch\_id} trata adequadamente a dependência intra-recorte, mas não torna o conjunto equivalente a uma amostra aleatória simples do espaço completo de cenas de sensoriamento remoto.
Os intervalos de confiança devem ser interpretados como medida de estabilidade dentro do desenho experimental adotado, e não como garantia automática de generalização irrestrita.

\subsection{Validade Externa}

A validade externa do bloco clássico é condicionada pelo conjunto de sensores e pelo protocolo de degradação artificial utilizados.
Os resultados são informativos para Sentinel-2, Landsat-8, Landsat-9 e MODIS sob máscaras e níveis de ruído controlados, mas não cobrem automaticamente outros sensores, outras resoluções espaciais, outros padrões de lacuna ou ruídos com distribuição distinta.

No bloco de aprendizado profundo, a limitação externa é ainda mais clara.
As arquiteturas foram treinadas e avaliadas em domínio Sentinel-2, de modo que as conclusões sobre U-Net, GAN, ViT, AE e VAE devem ser entendidas como válidas para esse contexto experimental específico.
Generalizar a mesma hierarquia arquitetural para cenário multissensor, para outras estações do ano ou para distribuições geográficas muito diferentes exigirá reavaliação empírica dedicada.

Por fim, o artigo analisa recortes estáticos e não incorpora dependência temporal explícita.
Em aplicações reais, lacunas podem ocorrer em sequências temporais nas quais redundância entre datas sucessivas altera substancialmente o problema.
Assim, as conclusões aqui obtidas são mais fortes para reconstrução espacial monodata do que para cenários completos de fusão espaço-temporal.

Em síntese, as ameaças à validade não anulam os achados centrais do estudo, mas definem seu escopo correto.
Os resultados sustentam com boa consistência comparações internas entre métodos clássicos e comparações internas entre arquiteturas de aprendizado profundo no desenho experimental adotado.
O que eles não sustentam, sem extensões adicionais, é uma afirmação universal sobre superioridade entre famílias metodológicas em qualquer sensor, qualquer regime de ruído e qualquer tarefa posterior.
